\documentclass[10pt,a4paper]{amsart}
\usepackage{amsmath,amssymb,amsthm,graphicx}
\usepackage{verbatim}
\newenvironment{reference}{\comment}{\endcomment}
%\newenvironment{reference}{}{}
\newenvironment{slogan}{\comment}{\endcomment}
\newenvironment{history}{\comment}{\endcomment}
% For commutative diagrams we use Xy-pic
\usepackage[all]{xy}

% We use 2cell for 2-commutative diagrams.
\xyoption{2cell}
% Theorem environments.
%
\theoremstyle{plain}
\newtheorem{theorem}[subsection]{Theorem}
\newtheorem{proposition}[subsection]{Proposition}
\newtheorem{lemma}[subsection]{Lemma}

\theoremstyle{definition}
\newtheorem{definition}[subsection]{Definition}
\newtheorem{example}[subsection]{Example}
\newtheorem{exercise}[subsection]{Exercise}
\newtheorem{situation}[subsection]{Situation}

\theoremstyle{remark}
\newtheorem{remark}[subsection]{Remark}
\newtheorem{remarks}[subsection]{Remarks}

\numberwithin{equation}{subsection}

% Macros
%
\def\lim{\mathop{\mathrm{lim}}\nolimits}
\def\colim{\mathop{\mathrm{colim}}\nolimits}
\def\Spec{\mathop{\mathrm{Spec}}}
\def\Hom{\mathop{\mathrm{Hom}}\nolimits}
\def\Ext{\mathop{\mathrm{Ext}}\nolimits}
\def\SheafHom{\mathop{\mathcal{H}\!\mathit{om}}\nolimits}
\def\SheafExt{\mathop{\mathcal{E}\!\mathit{xt}}\nolimits}
\def\Sch{\mathit{Sch}}
\def\Mor{\mathop{\mathrm{Mor}}\nolimits}
\def\Ob{\mathop{\mathrm{Ob}}\nolimits}
\def\Sh{\mathop{\mathit{Sh}}\nolimits}
\def\NL{\mathop{N\!L}\nolimits}
\def\CH{\mathop{\mathrm{CH}}\nolimits}
\def\proetale{{pro\text{-}\acute{e}tale}}
\def\etale{{\acute{e}tale}}
\def\QCoh{\mathit{QCoh}}
\def\Ker{\mathop{\mathrm{Ker}}}
\def\Im{\mathop{\mathrm{Im}}}
\def\Coker{\mathop{\mathrm{Coker}}}
\def\Coim{\mathop{\mathrm{Coim}}}

% Boxtimes
%
\DeclareMathSymbol{\boxtimes}{\mathbin}{AMSa}{"02}

%
% Macros for moduli stacks/spaces
%
\def\QCohstack{\mathcal{QC}\!\mathit{oh}}
\def\Cohstack{\mathcal{C}\!\mathit{oh}}
\def\Spacesstack{\mathcal{S}\!\mathit{paces}}
\def\Quotfunctor{\mathrm{Quot}}
\def\Hilbfunctor{\mathrm{Hilb}}
\def\Curvesstack{\mathcal{C}\!\mathit{urves}}
\def\Polarizedstack{\mathcal{P}\!\mathit{olarized}}
\def\Complexesstack{\mathcal{C}\!\mathit{omplexes}}
% \Pic is the operator that assigns to X its picard group, usage \Pic(X)
% \Picardstack_{X/B} denotes the Picard stack of X over B
% \Picardfunctor_{X/B} denotes the Picard functor of X over B
\def\Pic{\mathop{\mathrm{Pic}}\nolimits}
\def\Picardstack{\mathcal{P}\!\mathit{ic}}
\def\Picardfunctor{\mathrm{Pic}}
\def\Deformationcategory{\mathcal{D}\!\mathit{ef}}
\usepackage{cleveref}
\crefname{originalreference}{original reference}{original references}
\Crefname{originalreference}{Original reference}{Original references}
\makeatletter
\newlabel{morphisms-lemma-reach-points-scheme-theoretic-image}{{morphisms-lemma-reach-points-scheme-theoretic-image}{source}{}{sourcepointer-1}{}}
\newlabel{morphisms-lemma-reach-points-scheme-theoretic-image@cref}{{[originalreference][0][]morphisms-lemma-reach-points-scheme-theoretic-image}{[1][0][]source}{}{}{}}
\makeatother
\crefname{originalreference}{original reference}{original references}
\Crefname{originalreference}{Original reference}{Original references}
\makeatletter
\newlabel{stacks-properties-lemma-UR-quasi-compact-above-x}{{stacks-properties-lemma-UR-quasi-compact-above-x}{source}{}{sourcepointer-2}{}}
\newlabel{stacks-properties-lemma-UR-quasi-compact-above-x@cref}{{[originalreference][0][]stacks-properties-lemma-UR-quasi-compact-above-x}{[1][0][]source}{}{}{}}
\makeatother
\title{Specializations lift through a smooth presentation when the stack is quasi-compact above the point}
\author{The Stacks Project Authors}
\date{stacks-project:stacks-morphisms.tex:local-snapshot-20261012; GNU FDL 1.2 or later}
\begin{document}
\maketitle
\noindent\textit{Editorial note.} All delivered mathematical text is the author's own, copied from the pinned Stacks Project chapter source sources/stacks-algebraic-stacks-morphisms/pinned-original-source.tex; no statement or proof step was rewritten, reordered or removed.\par\medskip
\noindent\textit{Editorial note.} Preamble provenance: the Stacks Project keeps one single shared preamble for every chapter; this card takes it directly from the staged sources/stacks-algebraic-stacks-morphisms/preamble.tex as three verbatim spans (lines 8-12, 14-18 and 154-217), so no derived file was written and nothing was added or redefined. The wrapper class and the base packages are supplied by the assembler.\par\medskip
\noindent\textit{Editorial note.} Bibliography provenance: neither the counted statement nor its proof contains a citation (this lane rescans the whole span before writing the card), so the card carries no bibliography and no entry was reconstructed.\par\medskip
\noindent\textit{Editorial note.} Selection provenance: this card was opened under the file-based control contract (CONTROL.md), with the mechanical pre-card gate run on its statement and proof interval before the card was written.\par\medskip
\noindent\textit{Allowed standard prerequisites (editorial).}\begin{itemize}
\item The condition on a pair (stack, point) given by the quasi-compactness of the fibre above the point of the relation of a smooth presentation. Referred to in the statement (source lines 10341-10343) and used in the proof (source lines 10356-10358) as a pointer to the chapter on properties of stacks. Makes the fibre of the second projection over a chosen point a quasi-compact algebraic space.
\item The openness of the map on topological spaces induced by a smooth morphism, and the criterion for a point of a target to be reached by a point of a scheme mapping to it, through the scheme theoretic image. Both are used inside the counted proof (source lines 10354-10355 and 10369-10372), the second as a cross-chapter pointer. Produces a point of the presentation over the specialising point and then a point above the given specialisation specialising to the chosen point.
\end{itemize}\medskip
\ifdefined\hypertarget\hypertarget{sourcepointer-1}{}\fi
\noindent\textit{Source reference.} Original reference morphisms-lemma-reach-points-scheme-theoretic-image: sources/stacks-algebraic-stacks-morphisms/pinned-original-source.tex, source lines 2141--2141. Navigation pointer only: the counted text refers to the Stacks item morphisms-lemma-reach-points-scheme-theoretic-image, which belongs to another chapter of the project and is not part of this card. The pointer carries the target's own label as its number because this chapter's staged snapshot does not contain the other chapter's tag table; its content is not reproduced here.\par\medskip
\ifdefined\hypertarget\hypertarget{sourcepointer-2}{}\fi
\noindent\textit{Source reference.} Original reference stacks-properties-lemma-UR-quasi-compact-above-x: sources/stacks-algebraic-stacks-morphisms/pinned-original-source.tex, source lines 7518--7518. Navigation pointer only: the counted text refers to the Stacks item stacks-properties-lemma-UR-quasi-compact-above-x, which belongs to another chapter of the project and is not part of this card. The pointer carries the target's own label as its number because this chapter's staged snapshot does not contain the other chapter's tag table; its content is not reproduced here.\par\medskip
\begin{lemma}
\label{lemma-specialization}
Let $\mathcal{X}$ be an algebraic stack. Let $f : U \to \mathcal{X}$
be a smooth morphism where $U$ is an algebraic space. Let
$x' \leadsto x$ be a specialization of points of $|\mathcal{X}|$.
Let $u \in |U|$ with $f(u) = x$. If $(\mathcal{X}, x')$ satisfy the
equivalent conditions of Properties of Stacks,
Lemma \ref{stacks-properties-lemma-UR-quasi-compact-above-x},
then there exists a specialization
$u' \leadsto u$ in $|U|$ with $f(u') = x'$.
\end{lemma}

\begin{proof}
Choose an \'etale morphism $(U_1, u_1) \to (U, u)$ where $U_1$
is an affine scheme. Then we may and do replace $U$ by $U_1$.
Thus we may assume $U$ is an affine scheme.
Consider the algebraic space $R = U \times_\mathcal{X} U$
with smooth projections $t, s : R \to U$.
Choose a point $w \in U$ mapping to $x'$; this is possible as
$f : |U| \to |\mathcal{X}|$ is open.
By our assumption on $x'$ the fibre
$F' = t^{-1}(w) = R \times_{t, U} w$
of $t : R \to U$ over $w$ is a quasi-compact algebraic space.
Choose an affine scheme $T$ and a surjective \'etale morphism
$T \to F'$. The fact that $x' \leadsto x$ means that $u$ is in
the closure of the image of the morphism
$$
T \to F' \to R \xrightarrow{s} U
$$
Namely, this image is the fibre of $|U| \to |\mathcal{X}'|$
over $x'$; if some $u \in V \subset |U|$ open is disjoint
from this fibre, then $f(V)$ is an open neighbourhood of $x$
not containing $x'$; contradiction.
Thus by Morphisms, Lemma
\ref{morphisms-lemma-reach-points-scheme-theoretic-image}
we see that  there exists $u' \in |U|$ in the fibre
of $|U| \to |\mathcal{X}|$ over $x'$ which specializes to $u$.
\end{proof}

\end{document}
